% B_max 分析
\section{实验二：真正的物理瓶颈是什么？}\label{sec:bmax}

\textbf{论证目标}：展示框架的区分力——即使拓扑性能相同（$t^* = 1.0$），物理天花板在不同设计点上有数量级的差异。

\textbf{论证逻辑}：
如果只问"800 Gbps下是否feasible"，所有Dragonfly配置返回相同的答案。
但这是初筛工具的失败——它应该能区分设计点的优劣。
因此我们换一个问法：\textbf{不是"在800G下是否可行"，而是"最多能撑到多高带宽才被物理约束杀死"。}

\textbf{实验设计}：
对9个拓扑求解Valiant LP，取最优$\mathbf{L}^*$，代入Eq.~\ref{eq:bmax}计算$B_{\max}^{\text{geom}}$和$B_{\max}^{\text{thermal}}$。
绘制$B_{\max}$--$N$曲线，标注先绑定约束。

\textbf{预期发现}：
\begin{enumerate}
    \item \textbf{Bump预算始终先触顶}。$B_{\max}^{\text{geom}} < B_{\max}^{\text{thermal}}$在所有9个拓扑上成立，
    散热余量2--7倍。这意味着信号bump先用完——优先投资bump pitch缩减而非更激进的冷却。
    \item \textbf{$B_{\max} \propto 1/N$}。从$N=2$的391K Gbps到$N=24$的67K Gbps。
    可外推：$N=96$时$B_{\max} \sim 17$K Gbps。
    \item \textbf{在$N \le 30$和800 Gbps目标下，物理约束有80倍以上余量}。
    架构师可以安全地先专注于拓扑优化，物理验证后置。
\end{enumerate}

\textbf{这个实验的核心价值}：它展示了统一LP框架的区分力——不是"这个设计可不可行"，而是"这个设计的极限在哪里"。
这一信息直接指导技术路线规划。

\begin{figure}[!htbp]
    \centering
    \includegraphics[width=0.85\textwidth]{bmax.pdf}
    \caption{$B_{\max}$随$N$的变化。实线为bump预算上限$B_{\max}^{\text{geom}}$，虚线为散热上限$B_{\max}^{\text{thermal}}$。阴影区域为散热余量。}
    \label{fig:bmax}
\end{figure}

\begin{table}[!htbp]
    \centering
    \caption{B$_{\max}$分析——9个拓扑的物理带宽天花板}
    \label{tab:bmax}
    \footnotesize
    \begin{tabular}{lrrrl}
        \hline
        拓扑 & $N$ & $B_{\max}^{\text{geom}}$ (Gbps) & $B_{\max}^{\text{thermal}}$ (Gbps) & 先绑定约束 \\
        \hline
        DF(1,1,1) & 2  & 391,544 & 1,830,400 & geometry \\
        DF(2,1,1) & 6  & 161,453 & 468,523  & geometry \\
        DF(2,2,1) & 12 & 112,766 & 306,923  & geometry \\
        DF(2,3,1) & 18 & 87,718  & 229,895  & geometry \\
        DF(3,2,1) & 24 & 66,779  & 133,916  & geometry \\
        DF(2,4,1) & 24 & 71,389  & 183,405  & geometry \\
        Mesh(4)    & 16 & 66,925  & 469,295  & geometry \\
        Torus(4)   & 16 & 97,886  & 686,400  & geometry \\
        \hline
    \end{tabular}
\end{table}
